\section{同行、融资与组织复杂度}

前面几章讨论的是技术和产品链条，这一章补上创业公司的外部坐标。高继扬在访谈里多次提到同行：学习宇树的整机和供应链，学习 Physical Intelligence 的基础模型和人才密度，也观察智元等公司在管理、知识产权和组织上的动作。这种表达很有意思：他不是把同行只当竞争对手，而是把行业当成可学习系统。

\subsection{向同行学习什么}

本节把“学习同行”拆成更具体的能力来源：硬件公司提供供应链样本，模型公司提供智能样本，成熟管理团队提供组织样本。这样看，竞争对手也是行业知识的来源。

对宇树，他关注的是深入供应链：齿轮、电机、壳体、电磁仿真等底层能力。对 Physical Intelligence，他关注的是智能方向的领头羊位置、人才密度和资金密度。对智元，他关注的是成熟管理团队如何经营具身智能公司，包括知识产权、组织调整和实事求是的管理动作。

\begin{knowledgebox}{同行学习的三个维度}
\begin{tabularx}{\textwidth}{p{0.22\textwidth}p{0.32\textwidth}X}
\toprule
对象 & 可学习之处 & 对星海图的意义 \\
\midrule
宇树 & 整机、供应链、零部件自研和制造深度 & 补足机器人公司的物理底盘。 \\
Physical Intelligence & 基础模型、人才密度和前沿算法投入 & 对标机器人大脑的上游能力。 \\
智元 & 管理团队、知识产权、组织调整和经营动作 & 学习复杂组织如何持续交付。 \\
\bottomrule
\end{tabularx}
\end{knowledgebox}

\subsection{估值上涨不是组织问题的根源}

访谈里提到，星海图两年估值大幅上涨，团队从十几人扩到两百多人。高继扬认为，真正造成组织问题的不是估值上涨本身，而是组织变复杂、任务 scope 急速扩大、原有成员和创始团队能否跟上成长速度，以及能否及时引入更有经验的人。

具身智能公司还有一个特殊组织难题：整机供应链强调流程、纪律和质量；智能团队强调人才密度、创新和快速试错。这两种 domain 的管理语言不同，节奏不同，容错方式也不同。把它们放在同一家公司里，是具身智能创业天然的组织挑战。

\begin{importantbox}{具身智能公司的组织双重性}
一边是硬件/供应链，需要纪律、流程、质量和成本控制；一边是智能/模型，需要高密度人才、探索和快速迭代。能否同时管理这两种文化，是机器人公司能否长大的核心问题。
\end{importantbox}

\subsection{本章小结}

星海图不是在真空中竞争。它一边从同行学习，一边处理融资、组织扩张和双 domain 管理问题。对具身智能公司而言，技术路线能否兑现，最后会落到组织能否承受复杂度。

